% !TEX root = ../main.tex
\chapter{Background e Stato dell'Arte}\label{chap:background}

% PATTERN DISPONIBILI IN QUESTO CAPITOLO:
% - figure[ht]: \begin{figure}[ht]\centering\includegraphics[width=0.8\textwidth]{...}\caption{...}\label{...}\end{figure}
% - exmp: \begin{exmp} testo \end{exmp}  — esempi formali numerati per sezione
% - enumerate[label=(\alph*)]: per step/proprietà con etichette a), b)...
% - \footnote{\url{https://...}} — URL in nota a piè di pagina
% - \nameref{label} — cross-reference a una sezione per nome, non numero

\chapter{Background e Stato dell'Arte}\label{chap:background}

Il problema affrontato in questa tesi non è sorto nel vuoto: è il prodotto di una traiettoria evolutiva precisa, in cui la distribuzione della logica applicativa su reti di servizi autonomi ha moltiplicato la superficie esposta e reso obsoleti gli strumenti pensati per architetture monolitiche. Questo capitolo ricostruisce quella traiettoria in quattro movimenti. La Sezione~\ref{sec:paradigmi} analizza i paradigmi architetturali e i protocolli di comunicazione che definiscono il perimetro delle \gls{API} moderne. La Sezione~\ref{sec:gateway} esamina il pattern dell'\gls{API} gateway come punto di enforcement centralizzato della sicurezza nel paradigma Zero-Trust. La Sezione~\ref{sec:minacce} classifica le categorie di minaccia rilevanti per le \gls{API} REST esposte tramite gateway. La Sezione~\ref{sec:testing-evolution} analizza l'evoluzione degli approcci di testing, identificando i gap strutturali che motivano l'approccio contract-driven adottato nei capitoli successivi.


% ─────────────────────────────────────────────────────────────────────────────
\section{Paradigmi delle Architetture API Moderne}\label{sec:paradigmi}
% ─────────────────────────────────────────────────────────────────────────────

\subsection{Evoluzione dei Sistemi Distribuiti}\label{subsec:evoluzione-sistemi}

L'architettura monolitica, in cui tutta la logica applicativa risiede in un singolo processo deployabile, ha dominato lo sviluppo enterprise per decenni. Le sue proprietà sono note: deployment atomico, transazioni interne semplificate, toolchain di sviluppo consolidato. Il suo limite strutturale lo è altrettanto: la scala è ottenuta replicando l'intero monolite, indipendentemente da quale sottoinsieme funzionale sia effettivamente sotto carico, e ogni modifica richiede il redeployment dell'intera applicazione, aumentando il rischio di regressione con la crescita della codebase.

L'architettura a microservizi risponde a questo limite decompondo il sistema in servizi autonomi, ciascuno responsabile di un sottodominio delimitato del business, deployabile e scalabile indipendentemente dagli altri~\cite{MISSING:newman:2015:microservices}. % TODO: add to Zotero — Newman, S. (2015). Building Microservices. O'Reilly Media.
La decomposizione ridistribuisce la complessità: ogni servizio è più semplice del monolite, ma il sistema nel suo complesso introduce un nuovo ordine di problemi che il monolite non presentava. La comunicazione tra servizi avviene attraverso interfacce programmabili esposte in rete, ciascuna con il proprio contratto, il proprio ciclo di vita e la propria superficie di attacco. Il numero di queste interfacce cresce con le interazioni tra servizi, non linearmente con il numero di servizi in sé: un sistema a $n$ servizi può esporre fino a $O(n^2)$ canali di comunicazione, ciascuno potenzialmente raggiungibile, interrogabile e abusabile indipendentemente dal contesto applicativo.


\subsection{Tassonomia degli Stili Architetturali e di Comunicazione}\label{subsec:tassonomia-stili}

I servizi autonomi comunicano attraverso stili architetturali con proprietà tecniche distinte, ciascuno con implicazioni di sicurezza proprie. La rassegna che segue è orientata a quelle proprietà, non alla sintassi dei protocolli.

\gls{REST}, formalizzato da Fielding nella sua dissertazione del 2000~\cite{MISSING:fielding:2000:rest}, % TODO: add to Zotero — Fielding, R. T. (2000). Architectural Styles and the Design of Network-based Software Architectures. PhD dissertation, UCI.
è uno stile architetturale stateless basato su risorse identificate da \gls{URI} e manipolate tramite verbi \gls{HTTP} standard. La sua proprietà fondamentale è la separazione tra l'identificazione della risorsa e la sua rappresentazione: il client specifica quale risorsa vuole e quale operazione eseguire, il server decide come restituirla. Questa separazione genera una conseguenza di sicurezza strutturale: l'identificatore della risorsa è visibile nel path dell'\gls{URL} e controllabile dal client, rendendo la verifica dell'autorizzazione a livello di oggetto una responsabilità applicativa che il protocollo non può enforce autonomamente.

\gls{gRPC} adotta un modello diverso: le operazioni sono procedure nominate, i parametri sono tipizzati tramite Protocol Buffers, la serializzazione è binaria e il trasporto è \gls{HTTP}/2. La tipizzazione forte e la serializzazione binaria chiudono la superficie a certi vettori di injection, ma introducono opacità al livello di ispezione del payload: un \gls{WAF} che opera sulla semantica delle richieste \gls{HTTP}/1.1 è strutturalmente cieco ai frame Protobuf che attraversano un tunnel \gls{HTTP}/2.

\gls{GraphQL} espone un singolo endpoint che accetta query arbitrarie sul grafo dello schema: il client specifica esattamente i campi che vuole, compresi i percorsi di attraversamento del grafo. Questa flessibilità sposta il controllo dalla selezione degli endpoint alla validazione della complessità della query: una singola richiesta \gls{HTTP} può contenere centinaia di query alias, consentendo di aggirare rate limiting configurato a livello di conteggio delle richieste \gls{HTTP}.

\gls{SOAP}, protocollo basato su \gls{XML} e contratti \gls{WSDL}, è prevalente in sistemi enterprise legacy. Le sue vulnerabilità strutturali derivano dal parsing \gls{XML}: \gls{XXE} injection e XML Bomb sono conseguenze dirette del modello di elaborazione del documento, indipendenti dall'applicazione specifica.

WebSocket e \gls{SSE} condividono la proprietà di stabilire connessioni persistenti che sopravvivono al singolo ciclo request-response. Il rate limiting applicato sull'handshake \gls{HTTP} iniziale non si estende ai frame successivi, e i meccanismi di ispezione del gateway operano esclusivamente sul traffico di upgrade, lasciando il canale post-handshake fuori dal perimetro di enforcement.

Questa tesi si focalizza su \gls{API} \gls{REST} documentate tramite specifica OpenAPI, per le ragioni discusse nella Sezione~\ref{sec:gap-ricerca}: è lo stile architetturale più adottato nelle architetture a microservizi cloud-native e quello per cui esiste una specifica formale sfruttabile come oracolo di assessment.


\subsection{Il Fenomeno dell'API Sprawl}\label{subsec:api-sprawl}

La decomposizione in microservizi, combinata con cicli di sviluppo rapidi e team decentralizzati, produce un effetto collaterale documentato: la proliferazione incontrollata di endpoint esposti che sfugge a qualsiasi inventario centralizzato~\cite{MISSING:apisecurity:sprawl:2023}. % TODO: add to Zotero — ricerca web non ha restituito paper peer-reviewed specifico su API sprawl; le fonti disponibili sono report di settore (IBM, Gartner). Verificare esistenza di paper accademico prima di aggiungere.
Il fenomeno si manifesta in tre categorie distinte. Le \emph{Shadow API} sono endpoint funzionanti ma assenti da qualsiasi specifica pubblicata: creati durante lo sviluppo, mai rimossi, raggiungibili ma non documentati. Le \emph{Deprecated API} sono endpoint formalmente obsoleti che continuano a rispondere perché la loro disattivazione non è mai avvenuta o perché dipendenze non mappate li mantengono attivi. Le \emph{Zombie API} sono endpoint di versioni precedenti del servizio che persistono in deployment non aggiornati.

Tutte e tre le categorie condividono una proprietà critica per il security assessment: sono invisibili agli strumenti che costruiscono la propria conoscenza del target a partire dalla documentazione. Un fuzzer che enumera gli endpoint dalla specifica OpenAPI non può raggiungere un endpoint non documentato; uno scanner che confronta le risposte con i codici attesi dalla specifica non può rilevare una Shadow API perché non sa che dovrebbe cercarvela. La rilevazione delle Shadow API richiede un approccio inverso: confrontare la superficie effettivamente raggiungibile con quella dichiarata, identificando la discrepanza. Questo è precisamente uno degli obiettivi del Domain-0 (API Discovery) descritto nella Sezione~\ref{subsec:domini-tassonomia}.




% ─────────────────────────────────────────────────────────────────────────────
\section{Il Pattern dell'API Gateway in Ottica Zero-Trust}\label{sec:gateway}
% ─────────────────────────────────────────────────────────────────────────────

\subsection{Analisi Comparativa dei Modelli di Esposizione in Rete}\label{subsec:modelli-esposizione}

Le architetture a microservizi dispongono di quattro modelli principali per esporre i propri servizi al traffico esterno e per governare il traffico interno, ciascuno con un diverso posizionamento nel perimetro di sicurezza.

L'\gls{API} gateway enterprise è il punto di terminazione centralizzato per il traffico nord-sud, ovvero il traffico proveniente da client esterni. Concentra in un unico strato le funzioni trasversali che altrimenti andrebbero replicate in ogni servizio: autenticazione e autorizzazione, rate limiting, trasformazione dei payload, terminazione \gls{TLS} e logging delle transazioni. La centralizzazione rende il gateway il punto naturale in cui enforcare le policy di sicurezza prima che le richieste raggiungano la logica applicativa; è anche il punto in cui un errore di configurazione produce conseguenze sistemiche, non localizzate a un singolo servizio.

Il modello Kubernetes Ingress e Gateway \gls{API} governa il traffico nord-sud all'interno di un cluster containerizzato. La risorsa Ingress definisce regole di routing dichiarative (hostname, path, servizio di destinazione) che un Ingress Controller traduce in configurazione del reverse proxy sottostante. La Kubernetes Gateway \gls{API}, introdotta come standard stabile nel 2023, risolve le limitazioni di portabilità del modello Ingress sostituendo le annotation controller-specific con risorse tipizzate e role-oriented. Dal punto di vista del security assessment, il gateway gestisce il traffico di upgrade e l'autenticazione iniziale, ma la sua capacità di ispezione si limita al piano \gls{HTTP}: traffico WebSocket post-handshake e payload Protobuf attraversano il layer senza ispezione del contenuto.

Il service mesh governa il traffico est-ovest, ovvero le comunicazioni inter-servizio all'interno del perimetro. Architetture basate su sidecar proxy (Envoy nel caso di Istio) intercettano tutto il traffico del pod in modo trasparente, abilitando mutua autenticazione \gls{TLS} (mTLS), osservabilità e circuit breaking senza modifiche al codice applicativo. La mesh introduce però assunzioni implicite sulla natura del traffico: un sidecar che bufferizza messaggi per implementare retry e metriche assume che i flussi siano finiti e di durata limitata; uno stream gRPC bidirezionale indefinito viola questa assunzione e può produrre esaurimento di memoria nel proxy.

Il modello serverless delega la gestione del ciclo di vita dei container alla piattaforma cloud. Il traffico è gestito da un gateway managed che scala automaticamente e non richiede configurazione infrastrutturale esplicita. Il costo di questa astrazione è visibile sul piano della sicurezza: l'accesso amministrativo al piano di controllo è mediato da policy IAM della piattaforma e non è disponibile nella forma di un'Admin \gls{API} ispezionabile, il che rende inapplicabili le verifiche \texttt{WHITE\_BOX} descritte nella Sezione~\ref{subsec:box-gradient}.


\subsection{Criteri di Selezione e Giustificazione Architetturale}\label{subsec:criteri-selezione}

Dei quattro modelli analizzati, l'\gls{API} gateway enterprise è l'unico che soddisfa congiuntamente i requisiti del framework proposto in questa tesi. Tre criteri guidano questa selezione.

Il primo criterio è la verificabilità del piano di controllo. Un gateway enterprise espone un'Admin \gls{API} che consente l'ispezione diretta della configurazione: plugin attivi, policy di routing, parametri di rate limiting, variabili d'ambiente. Questa interfaccia è il presupposto tecnico dei test \texttt{WHITE\_BOX}, che verificano garanzie non osservabili dall'esterno tramite probe empirici. Il service mesh delega parte di questa visibilità al control plane distribuito; il modello serverless la sottrae interamente.

Il secondo criterio è la separazione tra piano di autenticazione e logica applicativa. Un gateway che centralizza l'enforcement delle policy di autenticazione e autorizzazione consente di valutare questi meccanismi indipendentemente dall'applicazione ospitata: un test che verifica il rifiuto delle richieste non autenticate non deve conoscere la logica applicativa, deve conoscere solo il contratto che il gateway dovrebbe enforcare. Questa separazione è la condizione che rende il framework applicazione-agnostico nel senso stabilito nella Sezione~\ref{subsec:agnosticismo}.

Il terzo criterio è la documentazione formale delle \gls{API} esposte. Un gateway enterprise instrada il traffico verso applicazioni che espongono la propria interfaccia tramite specifica OpenAPI; questa specifica è la sorgente da cui il framework deriva a runtime tutta la conoscenza del target. I modelli serverless e mesh non escludono questa possibilità, ma non la presuppongono come requisito architetturale.


\subsection{Enforcing del Paradigma Zero-Trust}\label{subsec:zero-trust}

Il paradigma Zero-Trust, formalizzato nel NIST Special Publication 800-207~\cite{MISSING:rose:2020:zerotrust}, % TODO: add to Zotero — Rose, S. et al. (2020). Zero Trust Architecture. NIST SP 800-207. DOI: 10.6028/NIST.SP.800-207
ridefinisce il perimetro di sicurezza spostando il punto di controllo dalla rete all'identità: nessuna richiesta è considerata attendibile per via della sua provenienza da una rete ritenuta fidata, e ogni accesso è valutato in modo indipendente rispetto all'identità del chiamante, al contesto della richiesta e ai permessi sulla risorsa specifica. Applicato alle \gls{API} REST, il principio si traduce in tre requisiti operativi: autenticazione obbligatoria su ogni endpoint, autorizzazione valutata per ogni coppia chiamante-risorsa, e assenza di fiducia implicita derivante dalla posizione di rete del client.

L'\gls{API} gateway è il punto architetturale in cui questi requisiti trovano enforcement naturale: si colloca tra il client e il backend, ha visibilità su ogni richiesta in transito, e può rifiutare le richieste che non soddisfano i controlli configurati prima che raggiungano la logica applicativa. La verifica che questo enforcement sia effettivamente attivo e correttamente configurato è esattamente il perimetro del framework proposto: le garanzie che il gateway dichiara nella propria configurazione (autenticazione richiesta su tutti gli endpoint, rate limiting attivo, header di sicurezza iniettati) devono corrispondere al comportamento osservabile a runtime, e questa corrispondenza non è verificabile senza un assessment sistematico. La distanza tra configurazione dichiarata e comportamento effettivo è il gap che motiva il dominio Configuration \& Hardening (Domain-6) e i test \texttt{WHITE\_BOX} della Sezione~\ref{subsec:box-gradient}.


% ─────────────────────────────────────────────────────────────────────────────
\section{Tassonomia delle Minacce e Vulnerabilità Logiche}\label{sec:minacce}
% ─────────────────────────────────────────────────────────────────────────────

\subsection{Il Semantic Gap e l'Evoluzione del Panorama delle Minacce}\label{subsec:semantic-gap}

La transizione verso architetture a microservizi ha spostato il baricentro delle vulnerabilità dalle anomalie sintattiche alle violazioni semantiche. Un \gls{WAF} tradizionale intercetta una SQL injection perché riconosce nel payload un pattern sintattico noto; la stessa infrastruttura è strutturalmente cieca a una richiesta \texttt{GET /users/456} perfettamente formata inviata da un utente che ha il permesso di leggere solo \texttt{/users/123}, perché quella violazione non si esprime nella forma della richiesta ma nella relazione tra l'identità del chiamante e la risorsa richiesta. Il \gls{WAF} ispeziona il protocollo; la vulnerabilità risiede nel contratto applicativo.

Questa distanza tra ciò che uno strumento di testing può osservare e ciò che una vulnerabilità semantica richiede di comprendere costituisce il \emph{semantic gap}: il divario strutturale tra la forma della transazione \gls{HTTP} e la semantica del contratto che quella transazione dovrebbe rispettare. Il framework \gls{OWASP} API Security Top~10 nella sua edizione 2023~\cite{MISSING:owasp:apisec:2023} % TODO: add to Zotero — OWASP Foundation (2023). OWASP API Security Top 10 2023. https://owasp.org/API-Security/editions/2023/en/0x11-t10/
formalizza questa transizione: le prime cinque categorie di rischio, Broken Object Level Authorization (API1), Broken Authentication (API2), Broken Object Property Level Authorization (API3), Unrestricted Resource Consumption (API4) e Broken Function Level Authorization (API5), sono tutte vulnerabilità semantiche. La loro manifestazione comune è una richiesta sintatticamente corretta che viola le assunzioni implicite del sistema sull'identità, i privilegi o il comportamento atteso del chiamante.

Colmare il semantic gap richiede che lo strumento di assessment possieda una rappresentazione formale del contratto \gls{API} e la utilizzi come oracolo per valutare le risposte ricevute. La specifica OpenAPI è quella rappresentazione: definisce quali endpoint esistono, quali parametri accettano, quali schemi di risposta producono e quali codici di stato sono attesi. Uno strumento che costruisce le proprie richieste nel rispetto di questa specifica e valuta le risposte rispetto a ciò che il contratto prevede per quella specifica operazione è in grado di raggiungere la classe di vulnerabilità che i fuzzer sintattici lasciano intatta. Questa è la premessa teorica dell'approccio adottato, il cui meccanismo implementativo è descritto nella Sezione~\ref{subsec:openapi-sot}.


\subsection{Categorie di Fallimento Architetturale}\label{subsec:categorie-fallimento}

Le vulnerabilità semantiche delle \gls{API} REST si raggruppano in tre macro-categorie distinte per natura del fallimento, in ordine crescente di complessità del vettore di attacco.

La prima categoria racchiude i difetti di autorizzazione e gestione dell'accesso agli oggetti. Il Broken Object Level Authorization (\gls{BOLA}, API1:2023) si manifesta quando un sistema verifica che il chiamante sia autenticato ma non verifica che abbia il diritto di accedere alla specifica risorsa identificata dal path. La richiesta è formalmente corretta; il fallimento risiede nell'assenza o nell'inadeguatezza del controllo sulla coppia chiamante-oggetto. Il Broken Function Level Authorization (\gls{BFLA}, API5:2023) segue la stessa logica a livello di operazione: un utente con un ruolo ristretto riesce a invocare un endpoint riservato a ruoli privilegiati, tipicamente perché i controlli di autorizzazione sono applicati in modo non uniforme tra endpoint dello stesso servizio. Entrambe le classi richiedono, per essere rilevate in modo sistematico, la capacità di effettuare richieste con credenziali di ruoli distinti e confrontarne le risposte: un workflow multi-step che i fuzzer stateless non supportano.

La seconda categoria comprende le violazioni del principio di minimo privilegio nel flusso dei dati. Il mass assignment si verifica quando un endpoint accetta nel payload di richiesta campi non documentati nello schema contrattuale, che il backend mappa direttamente sul modello di dominio senza whitelist esplicita: un attaccante può elevare i propri privilegi o modificare attributi riservati semplicemente includendo campi aggiuntivi in un \texttt{POST} o \texttt{PUT}. L'esposizione eccessiva di dati (API3:2023) è la controparte nella risposta: il servizio restituisce campi sensibili, tra cui hash di password, chiavi di \gls{API} o dati personali identificativi, perché il livello applicativo serializza l'intero oggetto di dominio affidando al client la selezione dei campi rilevanti. Entrambi i fallimenti sono invisibili a uno strumento che non conosce lo schema contrattuale atteso: a livello \gls{HTTP}, la transazione appare corretta in entrambi i sensi.

La terza categoria abbraccia l'abuso della topologia distribuita come vettore di attacco indiretto. Il \gls{SSRF} (API7:2023) sfrutta endpoint che accettano \gls{URL} forniti dal client, come webhook, mirror o proxy, per far sì che il server effettui richieste verso infrastrutture interne o verso i metadata endpoint dei cloud provider, esfiltrando credenziali di esecuzione o raggiungendo servizi non esposti al perimetro esterno. L'assenza di rate limiting (API4:2023) espone il sistema all'esaurimento controllato delle risorse: un attaccante che può inviare richieste senza limitazione può saturare il backend, i pool di connessioni al database o la capacità computazionale, rendendo il servizio indisponibile per gli utenti legittimi. Entrambe le classi operano attraverso funzionalità legittime del sistema, non attraverso payload malformati, e sono quindi sistematicamente invisibili ai layer di ispezione sintattica.


\subsection{Dal Modello Teorico alla Necessità di Valutazione Strutturata}\label{subsec:da-modello-a-valutazione}

Le tre categorie identificate nella sezione precedente non sono omogenee né per il tipo di precondizione che richiedono né per il meccanismo di rilevazione che presuppongono. Rilevare una violazione \gls{BOLA} richiede token validi per almeno due identità distinte e la capacità di confrontare le risposte ricevute con credenziali diverse. Rilevare mass assignment richiede di conoscere lo schema di richiesta contrattuale e di verificare il comportamento del sistema in presenza di campi non dichiarati. Rilevare un \gls{SSRF} richiede endpoint che accettino \gls{URL} e un meccanismo per osservare le richieste che il server effettua verso l'esterno. Nessuno dei tre fallimenti è raggiungibile con lo stesso strumento e la stessa configurazione.

Questa eterogeneità è la ragione per cui un singolo scanner generico non è sufficiente: le classi di vulnerabilità semantica impongono precondizioni di accesso distinte, richiedono strategie di osservazione diverse e non si mappano sulla stessa superficie di attacco. La risposta architetturale a questa eterogeneità è la tassonomia a otto domini di assurance descritta nella Sezione~\ref{subsec:domini-tassonomia}, in cui ogni dominio raggruppa le garanzie che condividono le stesse precondizioni informative e lo stesso tipo di evidenza necessario per la verifica. Il box gradient, discusso nella Sezione~\ref{subsec:box-gradient}, formalizza come il livello di visibilità disponibile determini la classe di garanzie raggiungibile: la struttura di precondizioni che emerge dall'analisi delle minacce in questa sezione non è un ostacolo da aggirare, ma il criterio che orienta la classificazione dei test nel Capitolo~\ref{chap:methodology}.


% ─────────────────────────────────────────────────────────────────────────────
\section{L'Evoluzione del Testing di Sicurezza: dall'Analisi Cieca ai Contratti}%
\label{sec:testing-evolution}
% ─────────────────────────────────────────────────────────────────────────────

\subsection{Paradigmi Tradizionali e i Limiti del DAST/SAST}\label{subsec:dast-sast}

Gli approcci classici al security testing delle applicazioni web si dividono in due famiglie distinte per punto di osservazione. Il \gls{SAST} analizza il codice sorgente o i binari compilati senza eseguire l'applicazione, identificando pattern noti di vulnerabilità a livello sintattico: buffer overflow, uso di funzioni non sicure, flussi di dati non validati verso sink pericolosi~\cite{MISSING:chess:2004:secure-programming}. % TODO: add to Zotero — verificare esistenza paper accademico su SAST per API; candidato: Chess, B., West, J. (2007). Secure Programming with Static Analysis. Addison-Wesley.
Il \gls{DAST} interagisce con l'applicazione in esecuzione attraverso la sua interfaccia esterna, inviando richieste e osservando le risposte alla ricerca di comportamenti anomali: codici di errore rivelatori, stack trace nelle risposte, variazioni nei tempi di risposta correlate a payload specifici.

Entrambi i paradigmi, applicati alle \gls{API} REST, incontrano un limite strutturale prima ancora di raggiungere la logica applicativa. Il \gls{DAST} genera payload a partire da euristiche sintattiche: variazioni di parametri, payload di injection, valori ai limiti degli intervalli attesi. Questi payload vengono costruiti senza conoscere il contratto che l'endpoint dichiara; di conseguenza, una quota sostanziale delle richieste viola i tipi, i formati o i vincoli enumerati nello schema e viene respinta con \texttt{HTTP~400} dal layer di validazione dell'input, prima che qualsiasi logica di business venga valutata. Questo fenomeno, denominato \emph{combinatorial wall}, è la barriera strutturale che impedisce al fuzzing cieco di raggiungere sistematicamente la classe di vulnerabilità semantica descritta nella Sezione~\ref{subsec:categorie-fallimento}: lo spazio di tutti i payload sintatticamente possibili è astronomico, ma lo spazio dei payload strutturalmente validi rispetto al contratto è definito e navigabile, e soltanto quest'ultimo è in grado di interrogare la logica applicativa. Il \gls{SAST}, dal canto suo, analizza il codice ma non può osservare il comportamento a runtime in presenza di identità autenticate con ruoli distinti: una violazione \gls{BOLA} emerge dall'interazione tra credenziali, routing del gateway e logica di autorizzazione, non da un pattern sintattico nel codice sorgente isolato.

Lo strumento \gls{DAST} più diffuso nel contesto delle \gls{API}, OWASP ZAP~\cite{MISSING:owasp:zap:2023}, % TODO: add to Zotero — OWASP Foundation. OWASP Zed Attack Proxy (ZAP). https://www.zaproxy.org
è progettato per operare su applicazioni web generiche e supporta \gls{API} REST tramite importazione della specifica OpenAPI. Anche in questa modalità, tuttavia, ZAP utilizza la specifica per enumerare gli endpoint e costruire richieste valide, ma valuta le risposte rispetto a euristiche generiche di anomalia, non rispetto a ciò che il contratto prevede per quella specifica operazione su quella specifica risorsa con quelle specifiche credenziali. La distinzione è sostanziale: un campo sensibile restituito in una risposta \texttt{200} è un finding solo se lo schema di risposta contrattuale non lo dichiarava atteso; uno scanner privo di quella conoscenza non può rilevarlo come anomalia.


\subsection{Asimmetria Informativa e il Box Gradient}\label{subsec:asimmetria-box-gradient}

La capacità di rilevare una vulnerabilità non dipende solo dalla sofisticazione dello strumento, ma dalla quantità di informazione che lo strumento possiede sul sistema sotto test prima di iniziare la valutazione. Questa asimmetria informativa è il principio che sottende la classificazione Black/Grey/White Box nel testing del software~\cite{MISSING:myers:2004:art-software-testing}: % TODO: add to Zotero — Myers, G. J., Sandler, C., Badgett, T. (2011). The Art of Software Testing. Wiley. (3rd ed.) — verificare edizione citabile
un tester che non ha accesso al codice sorgente né alla configurazione interna del sistema opera nel regime Black Box, limitando la propria osservazione al comportamento esterno; un tester con accesso completo alla configurazione interna opera nel regime White Box, potendo verificare per ispezione diretta proprietà non osservabili dall'esterno.

Applicata alle \gls{API} REST esposte tramite gateway, questa distinzione acquista una precisione tecnica che va oltre la semplice metafora della visibilità. Alcune garanzie di sicurezza sono fisicamente raggiungibili solo in determinati regimi di visibilità: verificare che il gateway rifiuti le richieste non autenticate non richiede nulla più dell'accesso di rete, perché il gateway applica quel controllo a qualsiasi interlocutore; simulare un attaccante esterno privo di credenziali è l'unico modo per ottenere evidenza non contaminata da privilegi impliciti. Verificare una violazione \gls{BOLA} richiede almeno due identità con ruoli distinti, entrambe autenticate: senza quello stato il gateway respinge la richiesta prima che la logica di autorizzazione venga valutata, e l'esito è indistinguibile da un corretto rifiuto di autenticazione. Verificare che le variabili d'ambiente del gateway non contengano credenziali in chiaro richiede accesso al piano di configurazione amministrativo: nessun probe empirico verso gli endpoint esposti può restituire quella informazione.

Ne deriva che il gradiente Black/Grey/White Box, nell'approccio adottato in questa tesi, non è una tassonomia teorica applicata retroattivamente ai test, ma il prodotto delle precondizioni concrete che ciascuna garanzia di sicurezza impone al tester. La sua formalizzazione come proprietà architetturale (descritta nella Sezione~\ref{subsec:box-gradient}) è la conseguenza diretta di questa analisi: assegnare a un test un livello di visibilità senza derivarlo dalle sue precondizioni produce una classificazione decorativa, non operativa.


\subsection{Contract-Driven Security e Approcci Ibridi}\label{subsec:contract-driven}

Il limite strutturale del fuzzing cieco, l'incapacità di costruire sistematicamente richieste valide rispetto al contratto senza possedere il contratto stesso, ha motivato una linea di ricerca che utilizza le specifiche formali delle \gls{API} come punto di partenza per la generazione dei test. RESTler~\cite{MISSING:atlidakis:2019:restler}, % TODO: add to Zotero — Atlidakis, V., Godefroid, P., Polishchuk, M. (2019). RESTler: Stateful REST API Fuzzing. ICSE 2019. DOI: 10.1109/ICSE.2019.00083
il primo fuzzer stateful per \gls{API} REST presentato all'ICSE 2019, analizza la specifica OpenAPI del servizio target e inferisce le dipendenze produttore-consumatore tra i tipi di richiesta: se un endpoint \texttt{POST /resources} produce un \texttt{resource\_id} che un endpoint \texttt{GET /resources/\{id\}} consuma come parametro di path, RESTler genera automaticamente la sequenza nell'ordine corretto. Questo approccio abbatte il combinatorial wall costruendo richieste strutturalmente valide a partire dal contratto e consentendo al fuzzer di raggiungere la logica applicativa che il fuzzing cieco non raggiunge.

RESTler dimostra la praticabilità tecnica del paradigma contract-driven, ma il suo obiettivo primario è la generazione di sequenze di test per trovare crash e comportamenti inattesi, non la verifica sistematica di garanzie di sicurezza specifiche. La specifica OpenAPI viene utilizzata per costruire le richieste, ma le risposte sono valutate rispetto a euristiche di anomalia generali, non rispetto a ciò che il contratto prevede per quella combinazione di operazione, risorsa e identità del chiamante. Ne deriva che RESTler è efficace nell'identificare difetti di robustezza e nell'esplorare lo spazio degli stati raggiungibili, ma non è progettato per verificare se una risposta \texttt{200} contenga campi che il contratto non prevede, se un token revocato venga ancora accettato, o se un endpoint di configurazione del gateway esponga credenziali non cifrate.

Il gap che rimane aperto è precisamente quello che questa tesi intende colmare: un framework che utilizzi la specifica OpenAPI non solo per costruire richieste valide ma come oracolo formale per valutare la correttezza semantica delle risposte, che operi in modo agnostico rispetto al target applicativo specifico, e che garantisca la riproducibilità deterministica dei risultati come requisito verificabile empiricamente. Il Capitolo~\ref{chap:methodology} presenta l'architettura che realizza questi tre obiettivi, a partire dai principi fondamentali che li rendono possibili congiuntamente.